\documentclass[12pt,addpoints]{exam}
\usepackage{amsmath,amssymb,amsthm}
\usepackage{geometry}
\geometry{a4paper, margin=1in}
\usepackage{listings}
\usepackage{graphicx}
\usepackage{tikz}

% \printanswers

\title{Mock Examination (Advanced Level): Advanced Algorithms}
\author{Course: Advanced Algorithms (EX01--EX15)}
\date{Time Allowed: 120 minutes \\ Total Points: 100}

\begin{document}

\maketitle

\begin{center}
\fbox{\fbox{\parbox{5.5in}{\centering
\textbf{Instructions to Candidates:}
\begin{itemize}
    \item This exam contains 6 questions.
    \item Answer all questions in English.
    \item Show all working out and proofs clearly.
    \item No calculators or external materials are permitted.
\end{itemize}
}}}
\end{center}

\vspace{0.2in}
\gradetable[v][questions]
\newpage

\begin{questions}

%=============================================================================
% QUESTION 1
%=============================================================================
\question[20] \textbf{Computational Complexity and co-NP}

\begin{parts}
\part[5] Let $A \subseteq \{0,1\}^*$ be a language. Prove that if $A \in \text{NP}$ and there exists a polynomial-time reduction from its complement $\overline{A}$ to $A$ (i.e., $\overline{A} \le_p A$), then $\text{NP} = \text{co-NP}$.

\begin{solution}
To prove $\text{NP} = \text{co-NP}$, we show containment in both directions:
\begin{enumerate}
    \item \textbf{Show $\text{NP} \subseteq \text{co-NP}$:}
    Let $L \in \text{NP}$. Since $A$ is $\text{NP}$-complete (or simply $\text{NP}$-hard), we have $L \le_p A$. 
    By taking the complement, this implies $\overline{L} \le_p \overline{A}$.
    We are given that $\overline{A} \le_p A$. Thus, by transitivity of polynomial-time reductions, we have:
    $$\overline{L} \le_p \overline{A} \le_p A \implies \overline{L} \le_p A$$
    Since $A \in \text{NP}$ and $\overline{L} \le_p A$, it follows that $\overline{L} \in \text{NP}$. 
    By definition of $\text{co-NP}$, since the complement of $L$ is in $\text{NP}$, $L$ must be in $\text{co-NP}$. Therefore, $\text{NP} \subseteq \text{co-NP}$.
    
    \item \textbf{Show $\text{co-NP} \subseteq \text{NP}$:}
    Let $L' \in \text{co-NP}$. By definition, $\overline{L'} \in \text{NP}$.
    Since $\overline{L'} \in \text{NP}$, we have $\overline{L'} \le_p A$, which implies $L' \le_p \overline{A}$.
    Since $\overline{A} \le_p A$, we have:
    $$L' \le_p \overline{A} \le_p A \implies L' \le_p A$$
    Since $A \in \text{NP}$, it follows that $L' \in \text{NP}$. Therefore, $\text{co-NP} \subseteq \text{NP}$.
\end{enumerate}
Combining these two containments gives $\text{NP} = \text{co-NP}$.
\end{solution}

\part[5] Prove that the problem $4\text{-COLOR}$ is $\text{NP}$-complete, assuming that $3\text{-COLOR}$ is $\text{NP}$-complete. Describe the reduction function and prove its correctness.

\begin{solution}
\begin{enumerate}
    \item \textbf{Membership in NP:} Given a 4-coloring assignment $c: V \to \{1,2,3,4\}$, we can verify in polynomial time $O(|V| + |E|)$ if $c(u) \neq c(v)$ for all $\{u,v\} \in E$. Thus, $4\text{-COLOR} \in \text{NP}$.
    \item \textbf{Reduction $3\text{-COLOR} \le_p 4\text{-COLOR}$:} Given a graph $G = (V,E)$, we construct a new graph $G' = (V', E')$ as follows:
    \begin{itemize}
        \item Add one new vertex $x$, so $V' = V \cup \{x\}$.
        \item Connect $x$ to every vertex in $V$, so $E' = E \cup \{\{v, x\} \mid v \in V\}$.
    \end{itemize}
    This construction takes $O(|V|)$ time, which is polynomial.
    \item \textbf{Proof of Equivalence:}
    \begin{itemize}
        \item \textbf{Forward Direction ($\implies$):} Suppose $G$ is 3-colorable with coloring $c: V \to \{1,2,3\}$. Define a coloring $c': V' \to \{1,2,3,4\}$ on $G'$ by setting $c'(v) = c(v)$ for $v \in V$, and $c'(x) = 4$. Since $c$ is a valid 3-coloring and $x$ has a unique color not used by any vertex in $V$, no adjacent vertices in $G'$ share a color. Thus, $G'$ is 4-colorable.
        \item \textbf{Backward Direction ($\impliedby$):} Suppose $G'$ is 4-colorable with coloring $c': V' \to \{1,2,3,4\}$. Without loss of generality, assume $c'(x) = 4$. Since $x$ is connected to every vertex in $V$, no vertex $v \in V$ can be colored 4. Thus, the restriction of $c'$ to $V$ must only use colors from $\{1,2,3\}$. This yields a valid 3-coloring of $G$.
    \end{itemize}
\end{enumerate}
Thus, $4\text{-COLOR}$ is $\text{NP}$-complete.
\end{solution}

\part[10] Let $\text{NEG-SOS}$ be a variation of the Sum of Subsets problem where the elements and target value can be negative integers:
$$\text{NEG-SOS} = \{\langle a_1, \dots, a_m, b \rangle \mid \exists I \subseteq \{1, \dots, m\} \text{ s.t. } \sum_{i \in I} a_i = b \text{ where } a_i, b \in \mathbb{Z}\}$$
Prove that $\text{NEG-SOS}$ is $\text{NP}$-complete by reducing the standard $\text{SOS}$ problem (where all $a_i, b \in \mathbb{N}$) to $\text{NEG-SOS}$. Discuss why standard certifications still hold for negative numbers.

\begin{solution}
\begin{enumerate}
    \item \textbf{Membership in NP:} The certificate is the set of indices $I$. The certifier computes the sum $\sum_{i \in I} a_i$ and compares it to $b$. The arithmetic additions of integers take polynomial time, regardless of sign. Thus, $\text{NEG-SOS} \in \text{NP}$.
    \item \textbf{Reduction $\text{SOS} \le_p \text{NEG-SOS}$:}
    The standard $\text{SOS}$ problem is defined for positive integers. Since $\mathbb{N} \subset \mathbb{Z}$, any instance of $\text{SOS}$ is already a valid instance of $\text{NEG-SOS}$. The reduction function is the identity function:
    $$f(\langle a_1, \dots, a_m, b \rangle) = \langle a_1, \dots, a_m, b \rangle$$
    This is trivially computable in $O(1)$ time.
    \item \textbf{Proof of Equivalence:}
    \begin{itemize}
        \item If $\langle a_1, \dots, a_m, b \rangle \in \text{SOS}$, there exists $I \subseteq \{1, \dots, m\}$ such that $\sum_{i \in I} a_i = b$. Since the elements and target are integers, the same index set $I$ satisfies the equation in $\text{NEG-SOS}$, so $f(x) \in \text{NEG-SOS}$.
        \item If $f(x) \in \text{NEG-SOS}$, there exists $I \subseteq \{1, \dots, m\}$ such that $\sum_{i \in I} a_i = b$. Since the input elements $a_i$ and target $b$ are positive, this equation holds exactly in $\text{SOS}$, so $x \in \text{SOS}$.
    \end{itemize}
\end{enumerate}
Thus, $\text{NEG-SOS}$ is $\text{NP}$-complete.
\end{solution}
\end{parts}

%=============================================================================
% QUESTION 2
%=============================================================================
\question[15] \textbf{Approximation Algorithms: Christofides' Algorithm}

The \textbf{Metric Traveling Salesperson Problem (Metric TSP)} is a variant of the TSP where the graph is complete and the edge weights satisfy the triangle inequality: $w(u,v) \le w(u,x) + w(x,v)$ for all $u, v, x \in V$.
The Christofides-Serdyukov algorithm is a famous $3/2$-approximation algorithm for Metric TSP. It improves upon the basic $2$-approximation (which simply duplicates all edges of a Minimum Spanning Tree) by adding a minimum-weight perfect matching on the set of odd-degree vertices of the MST. This results in a connected multigraph where every vertex has an even degree, which by Euler's Theorem guarantees the existence of an Eulerian tour. A Hamiltonian cycle (TSP tour) is then obtained by shortcutting the Eulerian tour.

\begin{parts}
\part[3] State the steps of the \textbf{Christofides-Serdyukov 3/2-approximation algorithm} for the Metric Traveling Salesperson Problem (metric TSP).

\begin{solution}
\begin{enumerate}
    \item \textbf{Compute MST:} Compute a Minimum Spanning Tree $T$ of the input complete graph $G$.
    \item \textbf{Odd-Degree Vertices:} Find the set $O$ of vertices in $T$ that have odd degrees.
    \item \textbf{Perfect Matching:} Compute a minimum-weight perfect matching $M$ on the complete subgraph induced by the vertices in $O$.
    \item \textbf{Eulerian Multigraph:} Combine the edges of $T$ and $M$ to form a multigraph $H = T \cup M$. Since every vertex in $H$ has an even degree, $H$ is Eulerian.
    \item \textbf{Eulerian Tour:} Find an Eulerian tour $C_{\text{Euler}}$ in $H$.
    \item \textbf{Shortcuts:} Obtain a Hamiltonian cycle $C_{\text{TSP}}$ by traversing $C_{\text{Euler}}$ and skipping already visited vertices using direct shortcuts.
\end{enumerate}
\end{solution}

\part[4] Prove that the number of odd-degree vertices in any tree (or graph) is even, thereby justifying why a perfect matching on $O$ always exists.

\begin{solution}
Let $G = (V,E)$ be any graph, and let $d(v)$ denote the degree of vertex $v$.
By the Handshaking Lemma, the sum of all vertex degrees is equal to twice the number of edges:
$$\sum_{v \in V} d(v) = 2|E|$$
Since the right-hand side $2|E|$ is an even integer, the sum on the left-hand side must also be even. 
We partition the sum into two groups: vertices of even degree and vertices of odd degree:
$$\sum_{v \in V \text{ s.t. } d(v) \text{ is even}} d(v) + \sum_{v \in O} d(v) = 2|E|$$
The first term is a sum of even numbers, which is always even. Thus, the second term $\sum_{v \in O} d(v)$ must also be even.
Since every term in the second sum $d(v)$ is odd, the only way for the sum of odd numbers to be even is if the number of terms $|O|$ is even. 
Since $|O|$ is even, we can partition $O$ into $|O|/2$ disjoint pairs, guaranteeing that a perfect matching on the induced subgraph of $O$ always exists.
\end{solution}

\part[8] Prove that the matching $M$ computed in step 3 satisfies $\text{cost}(M) \le \frac{1}{2} \text{OPT}$, and show how this completes the proof that the algorithm is a $3/2$-approximation.

\begin{solution}
\begin{enumerate}
    \item \textbf{Proof of $\text{cost}(M) \le \frac{1}{2} \text{OPT}$:}
    Let $C^*_{\text{TSP}}$ be the optimal TSP tour of $G$, so $\text{cost}(C^*_{\text{TSP}}) = \text{OPT}$.
    Consider the cycle $C_O$ obtained by traversing $C^*_{\text{TSP}}$ and skipping vertices not in $O$. By the triangle inequality (since the graph is metric), we have:
    $$\text{cost}(C_O) \le \text{cost}(C^*_{\text{TSP}}) = \text{OPT}$$
    The cycle $C_O$ has an even number of vertices $|O|$. We can partition the edges of $C_O$ into two disjoint perfect matchings on $O$, say $M_1$ and $M_2$.
    The sum of their costs is equal to the cost of $C_O$:
    $$\text{cost}(M_1) + \text{cost}(M_2) = \text{cost}(C_O) \le \text{OPT}$$
    Since the sum of the costs of the two matchings is at most $\text{OPT}$, the minimum-weight matching $M$ must satisfy:
    $$\text{cost}(M) \le \min(\text{cost}(M_1), \text{cost}(M_2)) \le \frac{1}{2} \text{OPT}$$
    
    \item \textbf{Proof of $3/2$-approximation:}
    The cost of the Eulerian tour $C_{\text{Euler}}$ is:
    $$\text{cost}(C_{\text{Euler}}) = \text{cost}(T) + \text{cost}(M)$$
    We know that $\text{cost}(T) \le \text{OPT}$ (removing an edge from the optimal TSP tour yields a spanning tree).
    Using the matching bound:
    $$\text{cost}(C_{\text{Euler}}) \le \text{OPT} + \frac{1}{2} \text{OPT} = \frac{3}{2} \text{OPT}$$
    By the triangle inequality, shortcutting does not increase the tour cost:
    $$\text{cost}(C_{\text{TSP}}) \le \text{cost}(C_{\text{Euler}}) \le \frac{3}{2} \text{OPT}$$
    This proves that Christofides' algorithm is a $3/2$-approximation.
\end{enumerate}
\end{solution}
\end{parts}

%=============================================================================
% QUESTION 3
%=============================================================================
\question[15] \textbf{Alpha-Beta Search Bounds and Proofs}

\begin{parts}
\part[7] In the best-case scenario (perfect move ordering), the number of leaf nodes evaluated by Alpha-Beta pruning in a complete game tree with branching factor $b$ and depth $d$ is given by the formula:
$$V(b,d) = b^{\lceil d/2 \rceil} + b^{\lfloor d/2 \rfloor} - 1$$
\begin{subparts}
    \subpart[4] Prove this formula for $d=3$. Explain which nodes must be evaluated (often called PV nodes or type 1 nodes) and which ones can trigger cutoffs.
    \subpart[3] Provide a brief intuitive explanation of why the exponents $\lceil d/2 \rceil$ and $\lfloor d/2 \rfloor$ represent the mathematical limit of minimax search reduction.
\end{subparts}

\begin{solution}
\begin{subparts}
    \subpart For $d=3$, the formula yields:
    $$V(b,3) = b^{\lceil 1.5 \rceil} + b^{\lfloor 1.5 \rfloor} - 1 = b^2 + b^1 - 1$$
    Let's analyze the leaf node evaluations in the game tree under optimal move ordering:
    \begin{itemize}
        \item \textbf{Evaluating the first child $L$ of the root (depth 1, Min):} 
        We must evaluate its first child $LL$ (depth 2, Max) completely, which requires searching all $b$ leaf children of $LL$ (Count: $b$ leaves). 
        Since $LL$ returns the optimal value for the Min node $L$, we obtain a beta bound $\beta = \text{val}(LL)$. For the remaining $b-1$ siblings of $LL$ under $L$, we only need to evaluate their first leaf child. Since move ordering is optimal, the first leaf child of each sibling will be worse (greater than or equal to $\beta$ for a Min node), immediately triggering a beta cutoff and pruning the remaining $b-1$ leaf children. This takes $1$ evaluation per sibling, totaling $b-1$ leaf evaluations.
        Thus, evaluating $L$ takes exactly $b + (b-1) = 2b-1$ leaf evaluations.
        
        \item \textbf{Evaluating the remaining $b-1$ children of the root, say $R$ (depth 1, Min):}
        Since $L$ has returned the minimax value $v$ of the root, we set the alpha bound at the root to $\alpha = v$. For each of the remaining $b-1$ siblings $R$ of $L$, under optimal move ordering, the value of $R$ is worse for Max (i.e., $\text{val}(R) \le \alpha$). 
        To prove $\text{val}(R) \le \alpha$ and trigger an alpha cutoff at the root, we only need to find one child of $R$ at depth 2 with value $\le \alpha$. Under optimal move ordering, the first child $RL$ (depth 2, Max) will be the best choice for Min node $R$, so it will satisfy $\text{val}(RL) \le \alpha$.
        Once we verify $\text{val}(RL) \le \alpha$, we trigger an alpha cutoff at the root for $R$, pruning the remaining $b-1$ children of $R$.
        However, to verify that the Max node $RL$ satisfies $\text{val}(RL) \le \alpha$, we must prove that ALL of its children (which are leaves) are $\le \alpha$. Thus, we cannot prune any children of $RL$; we must evaluate all $b$ of its leaf children.
        This requires $b$ leaf evaluations for each of the $b-1$ siblings $R$. The total number of leaf evaluations for these siblings is $(b-1) \times b = b^2 - b$.
        
        \item \textbf{Total Leaf Evaluations:}
        Summing these two components yields:
        $$\text{Total} = (2b - 1) + (b^2 - b) = b^2 + b - 1$$
        This exactly matches the formula $V(b,3) = b^{\lceil 3/2 \rceil} + b^{\lfloor 3/2 \rfloor} - 1 = b^2 + b - 1$.
    \end{itemize}
    
    \subpart To mathematically prove that a node in a minimax tree has a value of at least $v$:
    \begin{enumerate}
        \item The player whose turn it is (e.g., the Maximizer) only needs to find \textbf{one} move that leads to a value of at least $v$ (branching factor of $1$).
        \item However, to guarantee this value, they must verify that for \textbf{every possible reply} by the opponent (e.g., the Minimizer), the outcome remains at least $v$ (branching factor of $b$).
    \end{enumerate}
    As a result, down any search path of depth $d$, we only explore $1$ branch at our turns and must verify all $b$ branches at the opponent's turns. 
    This means there are $\lfloor d/2 \rfloor$ levels with branching factor $1$ (our choices) and $\lceil d/2 \rceil$ levels with branching factor $b$ (opponent's choices). The number of leaf evaluations is therefore dominated by $b^{\lceil d/2 \rceil} \cdot 1^{\lfloor d/2 \rfloor} = b^{\lceil d/2 \rceil}$.
    This represents the absolute mathematical limit of search reduction: we cannot prune any further because we must, at a minimum, verify one move for ourselves and all replies for the opponent.
\end{subparts}
\end{solution}

\part[8] Prove the \textbf{Fail-Low} and \textbf{Exact-Value} properties of Negamax Alpha-Beta search: If the search function \texttt{negamax(N, depth, alpha, beta)} returns a value $v$, then:
\begin{enumerate}
    \item If $v \le \alpha$, the true minimax score satisfies $S(N) \le v$.
    \item If $\alpha < v < \beta$, the true minimax score satisfies $S(N) = v$.
\end{enumerate}

\begin{solution}
We prove these properties by induction on the depth of the search.
\begin{itemize}
    \item \textbf{Base Case (depth = 0):} The node $N$ is a leaf. Negamax returns the static evaluation $v = S(N)$.
    \begin{enumerate}
        \item If $v \le \alpha$, then $S(N) = v \le v$, so $S(N) \le v$ holds.
        \item If $\alpha < v < \beta$, then $S(N) = v$ holds.
    \end{enumerate}
    
    \item \textbf{Inductive Step:} Assume both properties hold for all calls at depth $d-1$. Let $N$ be a node of depth $d$ with children $C_1, \dots, C_k$.
    Negamax iterates through the children, calling:
    $$v_i = -\texttt{negamax}(C_i, d-1, -\beta, -\alpha)$$
    Let $S(N) = \max_i (-S(C_i))$ be the true minimax score.
    \begin{enumerate}
        \item \textbf{Fail-Low ($v \le \alpha$):}
        If the algorithm returns $v \le \alpha$, it means that for all children $C_i$, the recursive call returned $v_i \le \alpha$.
        Thus, for all $i$:
        $$-v_i \le \alpha \implies \texttt{negamax}(C_i, d-1, -\beta, -\alpha) \ge -\alpha$$
        By the induction hypothesis (Fail-High case, since returned value is $\ge$ beta-bound of the child), we have:
        $$S(C_i) \ge -\alpha \implies -S(C_i) \le \alpha$$
        Since this holds for all children $C_i$:
        $$S(N) = \max_i (-S(C_i)) \le \alpha$$
        Since $v$ is the maximum value found (which is $\le \alpha$), it follows that $S(N) \le v$.
        
        \item \textbf{Exact Value ($\alpha < v < \beta$):}
        If the algorithm returns a value $v$ in the open interval $(\alpha, \beta)$, there must exist some child $j$ such that:
        $$v = -v_j \in (\alpha, \beta) \implies v_j \in (-\beta, -\alpha)$$
        Since $v_j$ falls strictly inside the child's window, by the induction hypothesis (Exact Value):
        $$S(C_j) = v_j \implies -S(C_j) = -v_j = v$$
        For all other children $i$, the search returned $v_i \ge -v$, which by induction implies:
        $$S(C_i) \ge v_i \implies -S(C_i) \le -v_i \le v$$
        Therefore:
        $$S(N) = \max( -S(C_j), \max_{i \neq j} (-S(C_i)) ) = \max(v, \le v) = v$$
        This proves $S(N) = v$.
    \end{enumerate}
\end{itemize}
\end{solution}
\end{parts}

%=============================================================================
% QUESTION 4
%=============================================================================
\question[15] \textbf{Connect Four Solver and Bitboards}

\begin{parts}
\part[5] In the $6 \times 7$ Connect Four bitboard representation, the shift value for checking anti-diagonal wins (top-left to bottom-right) is $S = 6$.
\begin{subparts}
    \subpart[2] Write the complete C/C++ code snippet to check for a 4-in-a-row anti-diagonal win.
    \subpart[3] Explain why $S = 6$ corresponds to this direction, referencing the row and column structure of the padded bitboard.
\end{subparts}

\begin{solution}
\begin{subparts}
    \subpart \textbf{Code Snippet:}
    \begin{lstlisting}[language=C]
uint64_t d2 = board & (board >> 6);
if (d2 & (d2 >> 12)) return true;
    \end{lstlisting}
    
    \subpart \textbf{Explanation:}
    Moving from one cell to its neighbor in the anti-diagonal direction means moving one column to the right (which increases the column index by 1) and one row down (which decreases the row index by 1).
    In our padded bitboard layout:
    \begin{itemize}
        \item Moving one column to the right increases the bit index by 7.
        \item Moving one row down decreases the bit index by 1.
    \end{itemize}
    Thus, the net change in bit index for an anti-diagonal step is $7 - 1 = 6$. A shift of 6 aligns every cell with its top-left neighbor.
\end{subparts}
\end{solution}

\part[10] Let \texttt{board} represent the current player's pieces, and \texttt{mask} represent the occupied board mask (containing pieces of both players). 
Write a C/C++ function \texttt{uint64\_t get\_winning\_moves(uint64\_t board, uint64\_t mask)} that returns a bitboard containing all legal winning moves for the current player on their next turn. 
A move is legal if and only if it is the lowest empty cell in its column (i.e., it sits directly above the top occupied cell in that column).

\begin{solution}
To find the winning moves, we must:
\begin{enumerate}
    \item Identify all empty cells that are legal moves (directly playable). 
    A playable cell in column $c$ is the cell immediately above the highest occupied cell in $c$. In bitboard representation, this corresponds to shifting the occupied \texttt{mask} up by 1 row (shift 1) and checking which cells are empty, plus adding row 0 cells if the column is completely empty:
    $$\texttt{legal\_moves} = (\texttt{mask} + \texttt{bottom\_mask}) \setminus \texttt{mask}$$
    where \texttt{bottom\_mask} contains 1s at row 0 of all columns. An easier way using addition:
    $$\texttt{legal\_moves} = (\texttt{mask} + \texttt{0x01020408102040}) \& \sim\texttt{mask}$$
    Wait, in Connect Four, since we can only drop pieces from the top, the playable cells can be found by adding the column bottom bits to the mask:
    $$\texttt{playable} = (\texttt{mask} \text{ shifted up by 1}) \cup \text{bottom row}$$
    And then subtracting the mask:
    $$\texttt{legal} = ((\texttt{mask} \ll 1) \lor \texttt{bottom}) \& \sim\texttt{mask}$$
    Here, \texttt{bottom} is \texttt{0x01020408102040} (since row 0 bits are 0, 7, 14, 21, 28, 35, 42, which correspond to $2^0, 2^7, 2^{14}, \dots$).
    
    \item For each of these legal cells, check if playing there creates a 4-in-a-row win. 
    We can test if a bit $p$ creates a win by temporarily setting it in `board` and running the win check, or we can use bitwise operations to check if the cell completes a 3-in-a-row.
\end{enumerate}

\textbf{C/C++ Implementation:}
\begin{lstlisting}[language=C]
#define BOTTOM 0x01020408102040ULL

// Helper function to check if a board has a win
bool is_winning_board(uint64_t b) {
    // Horizontal (shift 7)
    uint64_t h = b & (b >> 7);
    if (h & (h >> 14)) return true;
    // Diagonal (shift 8)
    uint64_t d1 = b & (b >> 8);
    if (d1 & (d1 >> 16)) return true;
    // Anti-diagonal (shift 6)
    uint64_t d2 = b & (b >> 6);
    if (d2 & (d2 >> 12)) return true;
    // Vertical (shift 1)
    uint64_t v = b & (b >> 1);
    if (v & (v >> 2)) return true;
    return false;
}

uint64_t get_winning_moves(uint64_t board, uint64_t mask) {
    // 1. Compute legal playable moves
    uint64_t legal = ((mask + BOTTOM) & ~mask) & 0x1F3E7CF9F3E7CFULL; // mask out sentinels
    uint64_t winning_moves = 0;
    
    // 2. Test each legal move bit
    uint64_t temp = legal;
    while (temp > 0) {
        uint64_t move = temp & -temp; // Extract lowest set bit (LSO)
        if (is_winning_board(board | move)) {
            winning_moves |= move;
        }
        temp &= temp - 1; // Clear lowest set bit
    }
    return winning_moves;
}
\end{lstlisting}
\end{solution}
\end{parts}

%=============================================================================
% QUESTION 5
%=============================================================================
\question[15] \textbf{Dynamic Programming: Bounded Levenshtein Distance}

\begin{parts}
\part[8] In spelling correction, we are often only interested in checking if the Levenshtein distance between two strings $X$ (length $m$) and $Y$ (length $n$) is at most $k$ (where $k \ll \min(m,n)$). 
Show how to optimize the classical $O(mn)$ dynamic programming algorithm to run in $O(k \cdot \min(m,n))$ time by only filling a diagonal band of width $2k+1$ around the main diagonal of the DP matrix. 
State the modified DP recurrence relation, including the boundary conditions for the band.

\begin{solution}
To compute the bounded Levenshtein distance:
\begin{enumerate}
    \item \textbf{DP Matrix Band:} For any entry $D(i, j)$, if $|i - j| > k$, the edit distance must exceed $k$ due to the difference in prefix lengths. Thus, we only need to compute $D(i, j)$ for indices satisfying:
    $$i - k \le j \le i + k$$
    This restricts the computation to a band of width $2k+1$ along the diagonal.
    \item \textbf{Boundary Conditions:}
    Any cell $D(i, j)$ lying outside the band is implicitly treated as having a value of $\infty$ (or $k+1$):
    $$D(i, j) = \infty \quad \text{if } |i - j| > k$$
    \item \textbf{Base Cases:}
    $$D(i, 0) = i \quad \text{for } 0 \le i \le k$$
    $$D(0, j) = j \quad \text{for } 0 \le j \le k$$
    Other base cells outside the band ($i > k$ or $j > k$) are initialized to $\infty$.
    \item \textbf{Recurrence Relation:} For $1 \le i \le m$ and $\max(1, i-k) \le j \le \min(n, i+k)$:
    $$D(i, j) = \min \begin{cases} D(i-1, j) + 1 & \text{if } j \ge i-k+1 \\ D(i, j-1) + 1 & \text{if } j \le i+k-1 \\ D(i-1, j-1) + \text{cost} & \text{always} \end{cases}$$
    where $\text{cost} = 0$ if $X[i] = Y[j]$, and $1$ otherwise.
\end{enumerate}
Since we only evaluate $2k+1$ cells per row, the total number of cells computed is at most $(2k+1) \cdot \min(m,n)$, yielding a running time of $O(k \cdot \min(m,n))$.
\end{solution}

\part[7] Explain how the concept of Aspiration Windows in Alpha-Beta search shares a similar logical philosophy with the Bounded Levenshtein Distance algorithm described in part (a). Focus on the handling of out-of-bounds results.

\begin{solution}
Both algorithms restrict computation by focusing on a narrow range of interest:
\begin{itemize}
    \item \textbf{Aspiration Windows} search within a narrow window $[\alpha, \beta]$ around an expected score. If the true value lies outside this window, the search returns a fail-low or fail-high bound, requiring a re-search with an expanded window.
    \item \textbf{Bounded Levenshtein Distance} restricts computation to a narrow diagonal band of width $2k+1$ corresponding to an edit distance threshold $k$. If the true distance is larger than $k$, the DP matrix values will propagate $\infty$, signifying that the edit distance lies outside our window of interest.
\end{itemize}
The key philosophy is \textbf{efficiency via optimism}: both algorithms assume the result lies in a narrow, favorable range. If it does, we save significant computation ($O(mn) \to O(kn)$ in DP, and massive branch cuts in game trees). If it does not, we quickly detect the boundary violation (``$\infty$'' in DP, fail-low/high in search) and adjust.
\end{solution}
\end{parts}

%=============================================================================
% QUESTION 6
%=============================================================================
\question[20] \textbf{RSA Cryptography: Optimization and Attacks}

\begin{parts}
\part[10] \textbf{RSA-CRT Optimization:} In practice, RSA decryption $m \equiv c^d \pmod n$ is accelerated using the Chinese Remainder Theorem (CRT) when the prime factors $p$ and $q$ of $n = pq$ are known.
\begin{subparts}
    \subpart[4] State the modular equations used to compute $m_p \equiv m \pmod p$ and $m_q \equiv m \pmod q$ directly from the ciphertext $c$ using the precomputed private key components $d_p = d \pmod{p-1}$ and $d_q = d \pmod{q-1}$. Explain why this is mathematically valid.
    \subpart[6] Show how to reconstruct the final message $m \pmod n$ using Garner's formula:
    $$m = m_q + q \cdot \left( q^{-1} \cdot (m_p - m_q) \pmod p \right)$$
    Explain why this CRT method is computationally faster than standard modular exponentiation $c^d \pmod n$.
\end{subparts}

\begin{solution}
\begin{subparts}
    \subpart \textbf{Equations:}
    $$m_p \equiv c^{d_p} \pmod p$$
    $$m_q \equiv c^{d_q} \pmod q$$
    \textbf{Mathematical Validity:}
    By Fermat's Little Theorem, for any prime $p$ and any integer $a \not\equiv 0 \pmod p$, we have $a^{p-1} \equiv 1 \pmod p$.
    Thus, we can reduce the exponent modulo $p-1$:
    $$m \equiv c^d \equiv c^{d \pmod{p-1}} \equiv c^{d_p} \pmod p$$
    This guarantees that the computed $m_p$ is equivalent to $m \pmod p$. The same logic holds for $m_q \pmod q$.
    
    \subpart \textbf{Garner's Formula Verification:}
    Let $h = q^{-1} \cdot (m_p - m_q) \pmod p$. Then $m = m_q + q \cdot h$.
    \begin{itemize}
        \item Modulo $q$: $m \equiv m_q + 0 \equiv m_q \pmod q$, which is correct.
        \item Modulo $p$: $m \equiv m_q + q \cdot q^{-1} (m_p - m_q) \equiv m_q + (m_p - m_q) \equiv m_p \pmod p$, which is correct.
    \end{itemize}
    Since $0 \le m_p < p$ and $0 \le m_q < q$, the value $m_q + q \cdot h$ lies in $[0, pq - 1]$, matching the range of the unique CRT solution.
    
    \textbf{Performance Benefit:}
    Modular exponentiation $x^y \pmod N$ has a complexity of $O(\log y \cdot \log^2 N)$ using schoolbook multiplication.
    - Standard RSA decryption uses a modulus $n$ of size $B$ bits.
    - RSA-CRT performs two exponentiations modulo $p$ and $q$, which are each of size $B/2$ bits.
    Since multiplication complexity is quadratic in bit-length, a modular multiplication modulo $p$ is $4 \times$ faster than modulo $n$.
    Moreover, the exponents $d_p$ and $d_q$ are only $B/2$ bits long (compared to the $B$-bit exponent $d$), reducing the number of multiplications by half.
    Thus, each of the two exponentiations takes $(1/4) \times (1/2) = 1/8$ of the work of a full exponentiation. 
    Running both takes $2 \times (1/8) = 1/4$ of the time of standard decryption, resulting in a $\approx 4\times$ speedup.
\end{subparts}
\end{solution}

\part[10] \textbf{Common Modulus Attack:} Suppose a sender encrypts the same message $m$ using two different public exponents $e_1$ and $e_2$ such that $\gcd(e_1, e_2) = 1$, using the same RSA modulus $n$.
An eavesdropper intercepts the two ciphertexts:
$$c_1 \equiv m^{e_1} \pmod n \quad \text{and} \quad c_2 \equiv m^{e_2} \pmod n$$
Show how the eavesdropper can recover the message $m$ in polynomial time without factoring $n$.

\begin{solution}
Since $\gcd(e_1, e_2) = 1$, by Bezout's Identity there exist integers $a$ and $b$ (which can be computed in polynomial time using the Extended Euclidean Algorithm) such that:
$$a \cdot e_1 + b \cdot e_2 = 1$$
Without loss of generality, assume that $a > 0$ and $b < 0$.
The eavesdropper can compute:
$$c_1^a \cdot (c_2^{-1})^{-b} \pmod n$$
where $c_2^{-1}$ is the modular multiplicative inverse of $c_2$ modulo $n$ (which exists and can be computed in polynomial time using the Extended Euclidean Algorithm since $\gcd(c_2, n) = 1$ under plain RSA).
Let's evaluate this expression:
$$c_1^a \cdot (c_2^{-1})^{-b} \equiv (m^{e_1})^a \cdot (m^{-e_2})^{-b} \equiv m^{a e_1} \cdot m^{b e_2} \equiv m^{a e_1 + b e_2} \equiv m^1 \equiv m \pmod n$$
Thus, the eavesdropper can recover $m$ by computing:
$$m \equiv c_1^a \cdot (c_2^{-1})^{-b} \pmod n$$
This takes $O(\log n)$ arithmetic operations (using Extended Euclidean Algorithm and binary exponentiation), which runs in polynomial time.
\end{solution}
\end{parts}

\end{questions}
\end{document}
